\section{Scope and open questions}\label{sec:limitations}

The t-exact realisation problem is settled by \Cref{cor:texact-classification}; projective-right-target kernel categories are classified up to dg Morita equivalence in \Cref{cor:projective-target-classes}. These conclusions do not classify arbitrary full-profile markings: \Cref{ex:marking-boundary} has a connective marking with nonzero \(H^{-2}\). Ordinary homological epimorphisms with infinite right target projective dimension, without t-exactness, and the general converse for Schur kernels remain outside the result.

The profile problem retains the marking. Coordinate symmetries already give abstract equivalences between distinct cuts (\Cref{ex:support-symmetry}); the open question is to classify such equivalences and identify invariants beyond these symmetries. We make no general Boolean-lattice assertion about the thick subcategories of an avoidance category. In particular, when \(Z=\varnothing\), the whole singularity category includes objects of full support; arguments using only closed-point supports for a nonempty excluded subset of \(\PP^1\) do not apply unchanged.

The support classification requires the pair-support criterion \ref{S:pair}. For arbitrary finite groups or self-injective algebras satisfying \textup{(Fg)}, that criterion must be checked or replaced by a suitable pair-support theory. The two-block example after \Cref{thm:pgroup} shows why individual supports alone are insufficient. The complete Nakayama profile lattices are computed only in ranks \(3,4\). The relation between higher multiplication and profile posets, and the possible dg algebras of non-silting presentation cones, remain the questions \Cref{q:profile-m3,q:nonsilting}.

The heart-amplitude converse over general noetherian rings needs the finite-generation condition \((\dagger_T)\). No general checkable replacement is supplied. Likewise, the multi-object finite-type dictionary for arbitrary infinite sets of localisation maps and the telescope property for all self-injective Nakayama algebras are not established here.

The negative-\(K\) summand obstruction vanishes in every example constructed here. Producing a localisation with nonzero obstruction remains \Cref{question:negative-k-counterexample}; the simultaneous-kernel formula is not an existence result. A further question is whether asymmetric eventual Ext-vanishing can produce two genuinely incomparable marked defects by explicit extension and coextension constructions. The phenomena in \cite{HunekeJorgensen03,JorgensenSega04} motivate this question, but no such localisation pair or witnesses to both failures of containment are provided here.

\subsection*{Statements and declarations}
\paragraph{Funding.} This research was self-funded by the author. No external funding was received.
\paragraph{Author contribution.} Marouane Lamharzi Alaoui is the sole author and is responsible for the mathematical arguments, manuscript, and computational supplement.
\paragraph{Code availability.} The exact verification scripts, captured outputs, and reproduction instructions are supplied in \texttt{anc/supplement/}; see \Cref{app:computation}. The computational files carry an MIT licence. No external datasets were used.

